\section{Introduction}\label{sec:intro}
\IEEEPARstart{D}{isassembly} supports the recovery of reusable components and materials from end-of-life products. A disassembly line balancing problem (DLBP) assigns selected tasks to ordered workstations. Task selection affects recovery value. Workstation allocation affects time, resource demand, and operating cost. Guo \etal\ describe the modeling and optimization challenges in disassembly sequence planning \cite{Guo2021}.

Human--robot collaboration provides several task execution options. A human can perform a manual task, a robot can perform an automated task, and a joint mode can use both resources. Guo \etal\ study stochastic operation time with a multi-objective shuffled frog leaping method \cite{Guo2024a}. Guo \etal\ also examine partial destructive disassembly with human--robot hybrid operation \cite{Guo2024}. Chu and Chen examine selective operations in a two-sided cooperative line \cite{Chu2025}. Huang \etal\ consider uncertain operation and equipment activation in a U-shaped line \cite{Huang2025}. These studies motivate explicit execution-mode and resource decisions. Their layouts and objectives differ from the multi-line setting considered here.

Coordination becomes important when several lines draw resources from common pools. Local plans can satisfy every task and station requirement, yet their combined staffing demand can exceed the available resources. Guo \etal\ model multiproduct collaborative hybrid lines with resource sharing and solve them with a discrete aquila optimizer \cite{Guo2025}. Gu \etal\ address shared transportation in multi-factory remanufacturing \cite{Gu2025}. Zhou \etal\ consider worker fatigue in multi-factory process optimization \cite{Zhou2025}. These studies establish the relevance of resource coordination and human factors. The present work focuses on fixed human and robot deployment across disassembly lines. It uses a station-level fatigue limit over one representative cycle. Its computational contribution concerns line-plan decomposition and auditable integer bounds, rather than the general introduction of resource sharing.

A compact mixed-integer linear programming (MILP) formulation provides a direct optimization model. Its relaxation can mix fractional task assignments within stations. An alternative formulation represents an entire feasible line plan by one column. All line-specific constraints then remain inside the pricing domain. Only two resource capacities connect the lines. Dantzig and Wolfe establish the decomposition principle for block-structured models \cite{Dantzig1960}. L\"ubbecke and Desrosiers explain the role of pricing and dual bounds in column generation (CG) \cite{Lubbecke2005}. This theory motivates the present formulation; decomposition itself is an established method.

The computational question is whether line-level convexification reduces the cost of obtaining profitable, certifiable integer plans. A stronger relaxation alone does not answer this question. Each pricing problem remains an integer optimization problem. Large line blocks can make pricing expensive. Fractional master solutions can also leave an integer gap. A practical framework therefore needs both line pricing and a separate mechanism for integer solutions and bounds.

This work makes three contributions. First, it formulates a profit-maximizing multi-line DLBP with selective tasks, three execution modes, cumulative fatigue, and shared resource pools. Second, it develops a whole-line reformulation with an integer equivalence proof, two-phase CG, and demand-driven resource-profile coordination. Third, it evaluates the framework against direct compact-model optimization on frozen synthetic inputs. The evaluation reports valid bounds and proof states as well as execution time. It also identifies a failure regime with few, large pricing blocks.

Section~\ref{sec:model} defines the problem and compact model. Section~\ref{sec:method} presents the decomposition and integer coordination. Section~\ref{sec:experiments} reports the computational evidence. Section~\ref{sec:conclusion} states the conclusions and remaining questions.
